\section{Conclusions}
\label{sec:conclusions}

Nuclear reactor physics conventionally incorporates microscopic quantum phenomena through reaction data while treating the internal nuclear state of the fuel as an effectively uncontrolled variable. This work has explored a generalization in which the nuclear state is promoted to an explicit dynamical variable,
\begin{equation}
\rho_N(t),
\end{equation}
and the effective fission response becomes
\begin{equation}
\sigma_f
=
\sigma_f(E,\rho_N).
\end{equation}

If an external actuator can modify the nuclear-state distribution,
\begin{equation}
u_Q
\rightarrow
\rho_N,
\end{equation}
then a complementary control pathway can be defined:
\begin{equation}
u_Q
\rightarrow
\rho_N
\rightarrow
\sigma_f^{\mathrm{eff}}
\rightarrow
R_f
\rightarrow
k_{\mathrm{eff}}.
\end{equation}

The individual physical ingredients of this chain are not entirely hypothetical. Nuclear quantum-optical experiments have demonstrated externally controlled nuclear excitation dynamics in selected systems \cite{Heeg2021,Yamaguchi2024}. More directly relevant to fission, measurements of the low-lying $^{235}$U isomer demonstrate that a change in nuclear state can produce a substantial change in neutron-induced fission cross section \cite{DEer1988}.

Thus, for a relevant fissile isotope, the relation
\begin{equation}
\rho_N
\rightarrow
\sigma_f
\end{equation}
is experimentally established.

What remains unresolved is the controllability of the first step,
\begin{equation}
u_Q
\rightarrow
\rho_N,
\end{equation}
at an efficiency, timescale, and addressed fraction large enough to matter.

Several mechanisms for populating $^{235m}$U have been investigated, including nucleus-electron coupling and inelastic electron scattering \cite{Chodash2016,LiuWang2022}, but no presently established method provides efficient and tunable control of a macroscopic isomeric population suitable for reactor applications.

The central physical question is therefore not simply whether fission probability depends on nuclear state. For $^{235}$U, it demonstrably does. The relevant questions are instead
\begin{equation}
\boxed{
\begin{minipage}{0.86\columnwidth}
\centering
Can the population of fission-distinct nuclear states be externally\\
and dynamically controlled?
\end{minipage}
}
\end{equation}
and, subsequently,
\begin{equation}
\boxed{
\begin{minipage}{0.86\columnwidth}
\centering
Can such microscopic control produce a non-negligible\\
macroscopic reactivity effect?
\end{minipage}
}
\end{equation}

The minimal two-state model shows how these questions can be expressed quantitatively. For the $^{235}$U ground-state/isomer pair,
\begin{equation}
\sigma_f^{\mathrm{eff}}
=
(1-p_m)\sigma_f^g
+
p_m\sigma_f^m,
\end{equation}
which provides a direct connection between a nuclear-state population and reactor quantities.

At the system level, the framework can be represented by an additional perturbative reactivity contribution,
\begin{equation}
\rho
=
\rho_{\mathrm{conv}}
+
\rho_Q,
\end{equation}
with
\begin{equation}
\rho_Q
=
\mathcal{G}[\rho_N].
\end{equation}

This formulation should be understood as a theoretical control framework, not as an immediately realizable reactor technology. State-preparation efficiency, the controllable fraction of nuclei, penetration and spatial uniformity of the actuator, relaxation and decoherence, competing reaction channels, energy requirements, and coupling to neutron kinetics remain major open problems.

The literature also constrains the novelty claim. State-dependent fission in $^{235m}$U is known, and prior work has already mentioned its possible use for fission control \cite{LiuWang2022}. The contribution proposed here is therefore narrower: the explicit formulation of nuclear-state population as a dynamical control variable coupled through a multiscale chain to fission rate, neutron multiplication, and reactor reactivity.

In this formulation, conventional and nuclear-state control operate on complementary parts of the fission process:
\begin{equation}
\boxed{
\begin{minipage}{0.82\columnwidth}
\centering
conventional control $\rightarrow$ control of the neutron environment
\end{minipage}
}
\end{equation}
whereas
\begin{equation}
\boxed{
\begin{minipage}{0.82\columnwidth}
\centering
nuclear-state control $\rightarrow$ control of the nuclear response of the fuel.
\end{minipage}
}
\end{equation}

The immediate next step is not the design of a nuclear-state-controlled reactor. It is the identification, calculation, and experimental study of systems in which an externally adjustable actuator produces a measurable change in a fission-distinct nuclear-state population and hence a measurable change in neutron-induced fission under an otherwise fixed neutron field.

The operational criterion is
\begin{equation}
\boxed{
\left.
\frac{\partial R_f}
{\partial u_Q}
\right|_{\Phi}
\neq0.
}
\end{equation}

A verified nonzero response would supply the microscopic parameters required to determine whether nuclear-state control can remain a laboratory-scale nuclear effect or become a meaningful new degree of freedom in the dynamics of multiplying fission systems.
